\section*{Hoofdstuk \ref{ch:g8:winds-global-circulation} --- Wind en de algemene circulatie}

\begin{solution}{exo:g8:winds-global-circulation:1}
(a) De noordoostpassaat, uit het noordoosten. (b) De \omterm{def:g8:winds-global-circulation:global-winds}{westenwinden}, uit het
noordwesten. (c) De \omterm{def:g8:winds-global-circulation:global-winds}{polaire oostenwinden}, uit het oosten (uit het
noordoosten). (d) De zuidoostpassaat, uit het zuidoosten.
\end{solution}

\begin{solution}{exo:g8:winds-global-circulation:2}
(a) Warme lucht stijgt op en koude lucht daalt. (b) Laag. (c) Van hoge druk
naar lage druk. (d) Naar rechts.
\end{solution}

\begin{solution}{exo:g8:winds-global-circulation:3}
(a) Niet waar: een \omterm{def:g8:winds-global-circulation:breezes}{zeewind} waait van de \omterm{def:g2:water-on-earth:ocean}{zee} naar het land (overdag).
(b) Waar. (c) Niet waar: ze waaien uit het oosten (uit het noordoosten of
het zuidoosten). (d) Waar.
\end{solution}

\begin{solution}{exo:g8:winds-global-circulation:4}
De tekening toont een kringloop: langs de grond een pijl van de \omterm{def:g2:water-on-earth:ocean}{zee} naar het
land; boven het land een pijl omhoog; op hoogte een pijl van het land terug
boven de \omterm{def:g2:water-on-earth:ocean}{zee}; boven de \omterm{def:g2:water-on-earth:ocean}{zee} een pijl omlaag. H staat op de grond onder de
dalende lucht (boven de \omterm{def:g2:water-on-earth:ocean}{zee}), L onder de stijgende lucht (boven het land).
\end{solution}

\begin{solution}{exo:g8:winds-global-circulation:5}
(a) 1.41 keer zoveel grond. (b) 50~\%. (c) 26~\%. (d) Bij de \omterm{def:g3:continents-and-oceans:poles}{polen} komt het
zonlicht schuin aan en wordt het over een groot oppervlak verspreid, zodat
elke vierkante meter weinig energie krijgt.
\end{solution}

\begin{solution}{exo:g8:winds-global-circulation:6}
(a) \omterm{def:g8:winds-global-circulation:global-winds}{Polaire oostenwinden} (\ang{70} tot \ang{60} NB), \omterm{def:g8:winds-global-circulation:global-winds}{westenwinden} (\ang{60}
tot \ang{30} NB), noordoostpassaat (\ang{30} NB tot de \omterm{def:g3:continents-and-oceans:equator}{evenaar}),
zuidoostpassaat (\omterm{def:g3:continents-and-oceans:equator}{evenaar} tot \ang{30} ZB), \omterm{def:g8:winds-global-circulation:global-winds}{westenwinden} (\ang{30} tot
\ang{60} ZB), \omterm{def:g8:winds-global-circulation:global-winds}{polaire oostenwinden} (\ang{60} tot \ang{70} ZB). (b) Bij de
\omterm{def:g3:continents-and-oceans:equator}{evenaar}, in de stiltegordel. (c) Rond \ang{30} NB en \ang{30} ZB, onder de
dalende lucht.
\end{solution}

\begin{solution}{exo:g8:winds-global-circulation:7}
(a) Het land ($31 > 22$): de wind waait van de \omterm{def:g2:water-on-earth:ocean}{zee} naar het land, een
\omterm{def:g8:winds-global-circulation:breezes}{zeewind}. (b) De \omterm{def:g2:water-on-earth:ocean}{zee} ($21 > 17$): de wind waait van het land naar de
\omterm{def:g2:water-on-earth:ocean}{zee}, een \omterm{def:g8:winds-global-circulation:breezes}{landwind}. (c) Het land: $31 - 17 = 14$ graden; de \omterm{def:g2:water-on-earth:ocean}{zee}: $22 -
21 = 1$ graad. (d) De zonnestralen verwarmen maar een dun laagje grond, dat
snel opwarmt en afkoelt; de \omterm{def:g2:water-on-earth:ocean}{zee} verspreidt haar warmte over een dikke laag
water, dus verandert haar temperatuur tussen dag en nacht nauwelijks.
\end{solution}

\begin{solution}{exo:g8:winds-global-circulation:8}
(a) Uit het noorden. (b) Naar rechts. (c) Uit het noordoosten (het is de
noordoostpassaat). (d) Zonder draaiing zou ze uit het zuiden waaien; de
draaiing buigt haar naar links af; ze waait uit het zuidoosten.
\end{solution}

\begin{solution}{exo:g8:winds-global-circulation:9}
Bij de \omterm{def:g3:continents-and-oceans:equator}{evenaar} worden de grond en de \omterm{def:g2:water-on-earth:ocean}{zee} sterk door de Zon \emph{verwarmd};
de warme, vochtige lucht \emph{stijgt} op; terwijl ze opstijgt,
\emph{koelt} ze af, en haar waterdamp wordt \emph{\omterm{def:g4:clouds-rain-snow:cloud}{wolken}}, die buien
en onweer geven. Rond \ang{25} NB \emph{daalt} de lucht die bij de \omterm{def:g3:continents-and-oceans:equator}{evenaar}
opsteeg; eronder is de druk aan de grond \emph{hoog}; de dalende lucht
warmt op, haar \omterm{def:g4:clouds-rain-snow:cloud}{wolken} drogen op en het gebied blijft \emph{droog}: een
\omterm{def:g5:weather-and-climate:desert}{woestijn}.
\end{solution}

\begin{solution}{exo:g8:winds-global-circulation:10}
(a) De \omterm{def:g8:winds-global-circulation:global-winds}{westenwinden}. Nee: ze waaien uit het westen, tegen een schip in dat
naar het westen vaart. (b) Om de gordel van de \omterm{def:g8:winds-global-circulation:global-winds}{passaatwinden} te bereiken,
ten zuiden van ongeveer \ang{30} NB. (c) De noordoostpassaat, tussen de
\omterm{def:g3:continents-and-oceans:equator}{evenaar} en ongeveer \ang{30} NB. (d) De \omterm{def:g8:winds-global-circulation:global-winds}{westenwinden}, ten noorden van
ongeveer \ang{30} NB (rond \ang{40} NB, de breedte van de thuishaven).
\end{solution}

\begin{solution}{exo:g8:winds-global-circulation:11}
(a) $1 \div 0.5 = 2$: de tabel geeft 2. (b) $1 \div 0.866 \approx
1.1547$, dat is \textbf{1.15}. (c) Op \ang{30}: $1000 \div 1.15 \approx
\qty{870}{W}$ op elke vierkante meter (precies $1000 \times 0.866 =
\qty{866}{W}$). Op \ang{60}: $1000 \div 2 = \qty{500}{W}$.
\end{solution}

\begin{solution}{exo:g8:winds-global-circulation:12}
(a) Drie kringlopen: lucht stijgt op bij de \omterm{def:g3:continents-and-oceans:equator}{evenaar} en daalt rond \ang{30};
stijgt op rond \ang{60} en daalt rond \ang{30}; stijgt op rond \ang{60} en
daalt bij de \omterm{def:g3:continents-and-oceans:poles}{pool}. (b) L bij de \omterm{def:g3:continents-and-oceans:equator}{evenaar} en rond \ang{60}; H rond \ang{30} en
bij de \omterm{def:g3:continents-and-oceans:poles}{pool}. (c) Van \ang{30} naar de \omterm{def:g3:continents-and-oceans:equator}{evenaar} gaat de lucht aan het
oppervlak naar het zuiden en wordt ze naar rechts afgebogen:
noordoostpassaat. Van \ang{30} naar \ang{60} gaat ze naar het noorden en
wordt ze naar rechts afgebogen: ze komt uit het zuidwesten. Van de \omterm{def:g3:continents-and-oceans:poles}{pool} naar
\ang{60} gaat ze naar het zuiden, naar rechts afgebogen: polaire
oostenwind. (d) In de middelste cel stroomt de lucht aan het oppervlak naar
de \omterm{def:g3:continents-and-oceans:poles}{pool}, dat wil zeggen naar het noorden; naar rechts afgebogen gaat ze
naar het noordoosten: dat zijn de \omterm{def:g8:winds-global-circulation:global-winds}{westenwinden}.
\end{solution}

\begin{solution}{pb:g8:winds-global-circulation:1}
\textbf{1.} $59.6 - 23.5 = 36.1$ lengtegraden.

\textbf{2.} $36.1 \times 108 \approx 3899$: ongeveer \qty{3900}{km}.

\textbf{3.} Ja: $3900 - 3896 = 4$~km, een heel klein verschil.

\textbf{4.} Ja: beide liggen dicht bij \ang{14} NB, tussen de \omterm{def:g3:continents-and-oceans:equator}{evenaar} en
ongeveer \ang{30} NB, de gordel van de noordoostpassaat.

\textbf{5.} Uit het noordoosten (de pijlen wijzen naar het
zuidwesten).

\textbf{6.} Vooruit, en een beetje opzij: de wind komt van achter het
schip, iets van rechts, en duwt het naar het westen.

\textbf{7.} Hij hoort bij de \omterm{def:g8:winds-global-circulation:circulation}{algemene circulatie}: lucht blijft rond
\ang{30} NB \omterm{def:g1:landscapes-around-us:landform}{dalen} en bij de \omterm{def:g3:continents-and-oceans:equator}{evenaar} opstijgen, dus blijft de lucht aan het
oppervlak naar de \omterm{def:g3:continents-and-oceans:equator}{evenaar} stromen, en de draaiing van de Aarde buigt haar
altijd naar rechts af.

\textbf{8.} De passaatwind zou tegen het schip in waaien. De kapitein moet
naar het noorden varen, voorbij ongeveer \ang{30} NB, om de \omterm{def:g8:winds-global-circulation:global-winds}{westenwinden} te
vinden.

\textbf{9.} $7 \times 24 = \qty{168}{km}$ per dag.

\textbf{10.} $3 \times 168 = \qty{504}{km}$.

\textbf{11.} $3896 \div 168 \approx 23.2$ dagen.

\textbf{12.} Ongeveer \textbf{23 dagen}: minder dan een maand.
\end{solution}
